\newcommand{\CC}{C\nolinebreak\hspace{-.05em}\raisebox{.4ex}{\small\bf +}\nolinebreak\hspace{-.10em}\raisebox{.4ex}{\small\bf +}\xspace}

\newcommand{\desc}[1]{\mbox{\footnotesize #1}}
\newcommand{\descsm}[1]{\mbox{\small #1}}

\newcommand\BNFkind[1]{\hbox{\rm #1}}
\newcommand\BNFcase[1]{\hbox{\rm\hskip1.5em #1}}
\newcommand\BNFfirst[2]{\hbox{\rm #1:\hskip0.75em #2}}
\newcommand\BNFother[2]{\hbox{\rm \hphantom{#1:}\hskip0.75em #2}}

\newcommand{\expty}[2]{(#1, #2)}

\newcommand{\class}{class\xspace}
\newcommand{\lbound}{l}%{\calL}
\newcommand{\ubound}{u}%{\calU}
\newcommand{\ind}{k}

\newcommand{\C}{\mathcal{C}}
\renewcommand{\listof}[1]{\ensuremath{\langle\emptyBecomesVeryThin{#1}\rangle}}
\newcommand{\listofb}[1]{\ensuremath{\langle\verythinmathspace\bar{#1}\verythinmathspace\rangle}}


\newcommand{\ANY}{\mbox{\it Any}}
\newcommand{\EXTENDS}{\mbox{ extends}}
\newcommand{\LOBR}{\llbracket}
\newcommand{\ROBR}{\rrbracket}
\newcommand{\OBR}[1]{\LOBR #1\ROBR}
\newcommand{\SSmbox}[1]{\mbox{\scriptsize #1}}
\newcommand{\PlaceHolder}[1]{$\langle$\textit{#1}$\rangle$}
\newcommand{\NEGQQ}{\mskip-24\thinmuskip}

\newcommand{\commentalign}{{\hskip-1em}}
\renewcommand{\algorithmiccomment}[1]{\quad{\large$\circledast$}\rm\ #1 }
\renewcommand{\algorithmicindent}{1em}%

\makeatletter
\newcommand{\slantverbatim}{\def\verbatim@font{\slshape\ttfamily\hyphenchar\font\m@ne\@noligs}}
\makeatother

%%Change rule names to appear in square brackets
\newcommand{\rncite}[1]{[\rn{#1}]}
\renewcommand{\inflabel}[1]{%
  \ifsuppressrulenames\else
    \def\lab{#1}%
    \ifx\lab\empty
      \relax
    \else
      \rncite{\lab}%
  \fi\fi
}
\typicallabel{}
\newdimen\tightdim
\newcommand{\tighten}[2]{{\def\andalso{\def\QUUX[####1/####2]{\tightdim=2em\multiply\tightdim by ####1\divide\tightdim by ####2}\QUUX[#1]\hskip\tightdim}#2}}

\newcommand{\failed}{\mbox{\textcolor{red}{\ding{55}}}}
\newcommand{\tocheck}[1]{\textcolor{red}{#1}}
\newcommand{\QUES}[1]{\textbf{\emph{\textcolor{red}{#1}}}}
% make a big red TODO label
\newcommand{\TODO}[1]{\textbf{\emph{\textcolor{red}{TODO}}}: \textsf{\footnotesize #1}}

\makeatletter
\def\overcirclerightarrow{\mathpalette\overcirclerightarrow@}
\def\overcirclerightarrow@#1#2{\vbox{\ialign{$##$\cr
     \noalign{\kern -0.4ex\nointerlineskip}
     \raise0.18ex\hbox{$\scriptstyle\circ$}\mkern-3mu#1{-}\mkern-6mu\cleaders\hbox{$#1\mkern-2mu{-}\mkern-2mu$}\hfill
     \mkern-6mu{\to}\cr
     \noalign{\kern -1\p@\nointerlineskip}
     \hfil#1#2\hfil\cr}}}
\makeatother

\newenvironment{proofcase}
               {\list{}{\divide\leftmargin by 2\itemindent=-\leftmargin}%
                \item\relax}
               {\hfill$\checkmark$\endlist}


\newcommand{\kwd}[1]{\mathtt{#1}}

\newcommand\textquotes[1]{``#1''}

\newdimen\infrulepadding  \infrulepadding=0.3em

\renewcommand{\typesetrule}[2]{%    extra padding to prevent infrule rule from being mistaken for rep rule
   \setrulebody{{\displaystyle
      {\genfrac{}{}{0.9pt}{}
               {\begin{array}{@{}c@{}}#1\end{array}}
               {\begin{array}{@{\hspace{\infrulepadding}}c@{\hspace{\infrulepadding}}}#2\end{array}}}}}}


%Failed experiment:
% \makeatletter
% \renewcommand{\typesetrule}[2]{%    extra padding to prevent infrule rule from being mistaken for rep rule
%    \setrulebody{{\displaystyle
%       {\genfrac{}{}{0.9pt}{}
% %               {\begin{array}{@{}>{\show\@preamble\tightenpremises}c<{\endtighten}@{}}#1\end{array}}
%                {\vbox{\let\\\@RULEarraycr\halign{\hfil\@RULEtightenpremises{${}##{}$}\hfil\cr#1\crcr}}}
%                {\begin{array}{@{\hspace{\infrulepadding}}c@{\hspace{\infrulepadding}}}#2\end{array}}}}}}
% \newcommand{\@RULEarraycr}{%
%   ${\ifnum0=`}$\fi\@ifstar\@RULExarraycr\@RULExarraycr}
% \newcommand{\@RULExarraycr}{\@ifnextchar[\@RULEargarraycr{\ifnum0=`{$\fi}${}\cr}}
% \newcommand{\@RULEargarraycr}[1]{%
%   \ifnum0=`{$\fi}${}\ifdim #1>\z@ \@RULExargarraycr{#1}\else
%    \@RULEyargarraycr{#1}\fi}
% \newcommand{\@RULExargarraycr}[1]{\@tempdima #1\advance\@tempdima \dp \@arstrutbox
%    \vrule \@height\z@ \@depth\@tempdima \@width\z@ \cr}
% \newcommand{\@RULEyargarraycr}[1]{\cr\noalign{\vskip #1}}
% \newcommand{\@RULEtightenpremises}[1]{{\message{CONSIDERING #1}\setbox0\hbox{#1}%
%     \ifdim \wd0 > \linewidth
%         \message{TIGHTENING}%
%         \hbox to \linewidth{\def\andalso{\hskip 1em \hfil}#1}%
%     \else \box0 \fi}}
% \makeatother


\makeatletter
\renewcommand\infrule{\@ifstar{\@infrulez}{\@infrulew}}
\def\@infrulew{\@ifnextchar[{\@infruley}{\@infrulex}} %]
\def\@infrulez{\@ifnextchar[{\@infruleq}{\@infrulep}} %]
\def\@infruleq[#1]#2#3{{\infrulepadding=0pt\@infruley[#1]{#2}{#3}}}
\def\@infrulep#1#2#3{{\infrulepadding=0pt\@infrule{#1}{#2}{#3}}}
% \@infrulex and \@infruley are defined by bcprules.sty
\makeatother

\newcommand\tightline[1]{\hbox to \linewidth{\def\andalso{{}$\hskip 1em plus 1fil${}}{${}#1{}$}}}

%for quoting judgments and such
\makeatletter
\newcommand\judgments[1]{\begin{center}$\typesetax{#1}$\end{center}}
\makeatother

\newcommand{\ifexpr}[3]{(\mathrel{\mathsf{if}}#1\mathrel{\mathsf{then}}#2\mathrel{\mathsf{else}}#3)}
\newcommand{\defneq}{\mathrel{\hbox{\tt =}}}
\newcommand{\var}{\textit{var}}
\newcommand{\ilk}[1]{\textit{ilk}\tuple{#1}}
%\newcommand{\Id}{\textit{Id}}
\newcommand{\Id}{\ensuremath{\iota}}
\newcommand{\emptyseq}{\bullet}
\newcommand{\whennonempty}[2]{{\def\nonemptytempA{#1}\def\nonemptytempB{}\ifx\nonemptytempA\nonemptytempB\else#2\fi}}
\newcommand{\turnstile}{\vdash}

\newcommand{\ms}{\preceq}
\renewcommand{\bar}{\rep}
\newcommand{\meet}{\wedge}
\newcommand{\existentialmeet}{\sqcap}
\newcommand{\calM}{\mathcal{M}}
\newcommand{\calP}{\mathcal{P}}
\newcommand{\calS}{\mathcal{S}}
\newcommand{\calB}{\mathcal{B}}
\newcommand{\calL}{\mathcal{L}}
\newcommand{\calU}{\mathcal{U}}
\newcommand{\calT}{\mathcal{T}}
\newcommand{\quoted}[1]{\begin{quote}#1\end{quote}}
\newcommand{\exc}{\mathrel{\lozenge}}
\newcommand{\existsexc}{\mathrel{\lozenge}}
%\newcommand{\nexc}{\mathrel{\hbox to 0pt{$\mskip -1.4mu\not$\hss}\lozenge}}
\newcommand{\nexc}{\mathrel{\hbox to 0pt{$\mskip 0mu\not$\hss}\lozenge}}
\newcommand{\smalllozenge}{\vcenter{\hbox{\scalebox{.5}{$\lozenge$}}}}
\newcommand{\normallozenge}{\vcenter{\hbox{$\lozenge$}}}

\newcommand{\altlozenge}{\ooalign{\hfil$\vcenter{\hbox{$\lozenge$}}$\hfil\cr\hfil$\cdot$\hfil}}
\newcommand{\bexc}{\mathrel{\altlozenge}}
\newcommand{\bexcm}{\mathrel{\altlozenge}_\textrm{m}}
\newcommand{\bnexc}{\mathrel{\hbox to 0pt{$\mskip -1.4mu\not$\hss}\altlozenge}}

\newcommand{\countof}[1]{\mathchoice{\hbox{\tt\#}}{\hbox{\tt\#}}{\hbox{\footnotesize\tt\#}}{\hbox{\tiny\tt\#}}\tuple{#1}}
\newcommand{\countofb}[1]{\mathchoice{\hbox{\tt\#}}{\hbox{\tt\#}}{\hbox{\footnotesize\tt\#}}{\hbox{\tiny\tt\#}}\tupleb{#1}}
\newcommand{\bigcountof}[1]{\mathchoice{\hbox{\tt\#}}{\hbox{\tt\#}}{\hbox{\footnotesize\tt\#}}{\hbox{\tiny\tt\#}}\bigtuple{#1}}
\newcommand{\bigcountofb}[1]{\mathchoice{\hbox{\tt\#}}{\hbox{\tt\#}}{\hbox{\footnotesize\tt\#}}{\hbox{\tiny\tt\#}}\bigtupleb{#1}}
\newcommand{\Bigcountof}[1]{\mathchoice{\hbox{\tt\#}}{\hbox{\tt\#}}{\hbox{\footnotesize\tt\#}}{\hbox{\tiny\tt\#}}\Bigtuple{#1}}
\newcommand{\Bigcountofb}[1]{\mathchoice{\hbox{\tt\#}}{\hbox{\tt\#}}{\hbox{\footnotesize\tt\#}}{\hbox{\tiny\tt\#}}\Bigtupleb{#1}}
\newcommand{\LRcountof}[1]{\mathchoice{\hbox{\tt\#}}{\hbox{\tt\#}}{\hbox{\footnotesize\tt\#}}{\hbox{\tiny\tt\#}}\LRtuple{#1}}
\newcommand{\LRcountofb}[1]{\mathchoice{\hbox{\tt\#}}{\hbox{\tt\#}}{\hbox{\footnotesize\tt\#}}{\hbox{\tiny\tt\#}}\LRtupleb{#1}}

\newcommand{\apply}{\hbox{\tt\char'100}}
\newcommand{\comprises}{\mathrel{\mathord{\equiv}\mathord{\bigcup}}}

\newcommand{\fresh}[1]{\textit{fresh}({#1})}
\newcommand{\freeVar}[1]{\textit{freeVars({#1})}}

\newcommand{\bigtypeunion}{\bigsqcup}
\newcommand{\bigtypeintersection}{\bigsqcap}
\newcommand{\transposedtypeunion}{\mathop{\bigtypeunion^{\mathsf{T}}}}
\newcommand{\transposedtypeintersection}{\mathop{\bigtypeintersection^{\mathsf{T}}}}
\newcommand{\cartesiantypeunion}[1]{\mathop{\bigtypeunion^{\raise0.3ex\hbox{$\scriptstyle\times$}}_{#1}}}
\newcommand{\cartesiantypeintersection}[1]{\mathop{\bigtypeintersection^{\raise0.3ex\hbox{$\scriptstyle\times$}}_{#1}}}

\newcommand{\compose}{\circ}
\newcommand{\funcompose}{\bullet}

\newcommand{\excr}{\triangleright}
\newcommand{\excl}{\triangleleft}
\newcommand{\excrx}{\excr_\textrm{x}}
\newcommand{\exclx}{\excl_\textrm{x}}
\newcommand{\excrc}{\excr_\textrm{c}}
\newcommand{\exclc}{\excl_\textrm{c}}
\newcommand{\excro}{\excr_\textrm{o}}
\newcommand{\exclo}{\excl_\textrm{o}}

\newcommand{\excx}{\exc_\textrm{x}}
\newcommand{\excc}{\exc_\textrm{c}}
\newcommand{\exco}{\exc_\textrm{o}}
\newcommand{\excm}{\exc_\textrm{m}}

\newcommand{\propop}{\ensuremath{\mathrel{\ast}}}

\newcommand{\properancestorsname}{\textit{properAncestors}}
\newcommand{\properancestors}[2][\Delta]{\properancestorsname(#1,#2)}
\newcommand{\ancestorsname}{\textit{ancestors}}
\newcommand{\ancestors}[2][\Delta]{\ancestorsname(#1,#2)}
\newcommand{\excluded}[2][\Delta]{\textit{excluded}(#1,#2)}
\newcommand{\comprised}[2][\Delta]{\textit{comprised}(#1,#2)}
\newcommand{\parents}[2][\Delta]{\textit{parents}(#1,#2)}

\newcommand{\extends}{\ensuremath{<\mskip-1mu\mathrel{\raise0.09ex\hbox{:}}}}
\newcommand{\subtypeof}{\ensuremath{<\mskip-1mu\mathrel{\raise0.09ex\hbox{:}}}}
\newcommand{\subtypeofeq}{\ensuremath{\leq\mskip-1mu\mathrel{\raise0.3ex\hbox{:}}}}
\newcommand{\nsubtypeof}{\ensuremath{\nless\mskip-1mu\mathrel{\raise0.09ex\hbox{:}}}}
\newcommand{\supertypeof}{\ensuremath{\mathrel{\raise0.09ex\hbox{:}}}\mskip-1mu>}
\newcommand{\supertypeofeq}{\ensuremath{\mathrel{\raise0.3ex\hbox{:}}}\mskip-1mu\geq}
\newcommand{\leinner}{\ensuremath{\lesssim}}
\newcommand{\quantsubtypeof}{\ensuremath{\sqsubseteq}}
\newcommand{\morespec}{\ensuremath{\sqsubseteq}}

\newcommand{\alphaequiv}{\ensuremath{\stackrel{\alpha}{\sim}}}
\newcommand{\cequiv}{\ensuremath{\sim}}

\newcommand{\arrowtype}[2]{\mbox{\ensuremath{(\barearrowtype{#1}{#2})}}}
\newcommand{\barearrowtype}[2]{\mbox{\ensuremath{{#1} \rightarrow {#2}}}}
\newcommand{\tuple}[1]{\ensuremath{(\emptyBecomesVeryThin{#1})}}
\newcommand{\tupleb}[1]{\ensuremath{(\bar{#1})}}
\newcommand{\bigtuple}[1]{\ensuremath{\big(\emptyBecomesVeryThin{#1}\big)}}
\newcommand{\bigtupleb}[1]{\ensuremath{\big(\bar{#1}\big)}}
\newcommand{\Bigtuple}[1]{\ensuremath{\Big(\emptyBecomesVeryThin{#1}\Big)}}
\newcommand{\Bigtupleb}[1]{\ensuremath{\Big(\bar{#1}\Big)}}
\newcommand{\LRtuple}[1]{\ensuremath{\left(\emptyBecomesVeryThin{#1}\right)}}
\newcommand{\LRtupleb}[1]{\ensuremath{\left(\bar{#1}\right)}}

\newcommand{\uniontype}[2]{\mbox{\ensuremath{(\bareuniontype{#1}{#2})}}}
\newcommand{\bigbiguniontype}[2]{\mbox{\ensuremath{\bigl(\bareuniontype{#1}{#2}\bigr)}}}
\newcommand{\LRuniontype}[2]{\mbox{\ensuremath{\left(\bareuniontype{#1}{#2}\right)}}}
\newcommand{\bareuniontype}[2]{\mbox{\ensuremath{{#1} \sqcup {#2}}}}
\newcommand{\intersectiontype}[2]{\mbox{\ensuremath{(\bareintersectiontype{#1}{#2})}}}
\newcommand{\bigbigintersectiontype}[2]{\mbox{\ensuremath{\bigl(\bareintersectiontype{#1}{#2}\bigr)}}}
\newcommand{\LRintersectiontype}[2]{\mbox{\ensuremath{\left(\bareintersectiontype{#1}{#2}\right)}}}
\newcommand{\bareintersectiontype}[2]{\mbox{\ensuremath{{#1} \sqcap {#2}}}}
\newcommand{\existentialmeettype}[2]{\mbox{\ensuremath{\left({#1} \existentialmeet {#2}\right)}}}
\newcommand{\naryuniontype}[1]{\bigtypeunion\listof{#1}}
\newcommand{\naryintersectiontype}[1]{\bigtypeintersection\listof{#1}}
\newcommand{\naryuniontypeb}[1]{\naryuniontype{\bar{#1}}}
\newcommand{\naryintersectiontypeb}[1]{\naryintersectiontype{\bar{#1}}}

\newcommand{\minimalname}{\textit{minimal}}
\newcommand{\minimal}[1]{\ensuremath{\minimalname(#1)}}
\newcommand{\hasminimalname}{\textit{hasMinimal}}
\newcommand{\hasminimal}[1]{\ensuremath{\hasminimalname(#1)}}

\newcommand{\allvisible}[2][\Delta]{\ensuremath{\textit{allVisible}(#1,#2)}}
\newcommand{\visible}[2][\Delta]{\ensuremath{\textit{visible}(#1,#2)}}
\newcommand{\visiblefunc}[2][\Delta]{\ensuremath{\textit{visible}(#1,#2)}}
\newcommand{\visiblemeth}[3][\Delta]{\ensuremath{\textit{visible}(#1,(#2,#3))}}
\newcommand{\defined}[3][\Delta]{\ensuremath{\textit{defined}(#1,#2,#3)}}
\newcommand{\name}[1]{\ensuremath{\textit{name}\tuple{#1}}}
\newcommand{\bigname}[1]{\ensuremath{\textit{name}\bigtuple{#1}}}
\newcommand{\FLAT}[1]{\ensuremath{\textit{flat}(#1)}}
\newcommand{\flipname}{\textit{flip}}
\newcommand{\flip}[1]{\ensuremath{\flipname(#1)}}
\newcommand{\distinct}[1]{\ensuremath{\mathit{distinct}(#1)}}
\newcommand{\methods}[2][\Delta]{\ensuremath{\textit{methods}(#1,#2)}}
\newcommand{\applicablename}{\ensuremath{\textit{applicable}}}
\newcommand{\applicable}[3][\Delta]{\ensuremath{\applicablename\tuple{#1,#2,#3}}}
\newcommand{\bigapplicable}[3][\Delta]{\ensuremath{\applicablename\bigtuple{#1,#2,#3}}}
\newcommand{\instantiate}[2]{\ensuremath{\textit{instantiate}\tuple{#1,#2}}}
\newcommand{\biginstantiate}[2]{\ensuremath{\textit{instantiate}\bigtuple{#1,#2}}}
\newcommand{\firstdispatchable}[4][\Delta]{\ensuremath{\textit{firstDispatchable}(#1,#2,#3,#4)}}
\newcommand{\firstapplicable}[3][\Delta]{\ensuremath{\textit{firstApplicable}(#1,#2,#3)}}
\newcommand{\mostspecific}[2][\Delta]{\ensuremath{\textit{mostSpecific}(#1,#2)}}
\newcommand{\msavname}{\textit{msav}}
\newcommand{\msav}[3][\Delta]{\ensuremath{\msavname(#1,#2,#3)}}

\newcommand{\topologicalsort}[2][\Delta]{\ensuremath{\textit{topologicalSort}(#1, #2)}}
\newcommand{\orderedin}[3]{\ensuremath{\Delta\,\turnstile\, #1\;\mathsf{and}\;#2\;\mathsf{ordered}\;\mathsf{in}\;#3}}

\newcommand{\abstractdecls}[1]{\ensuremath{\textit{abstractDecls}\tuple{#1}}}
\newcommand{\concretedecls}[1]{\ensuremath{\textit{concreteDecls}\tuple{#1}}}
\newcommand{\bigabstractdecls}[1]{\ensuremath{\textit{abstractDecls}\bigtuple{#1}}}
\newcommand{\bigconcretedecls}[1]{\ensuremath{\textit{concreteDecls}\bigtuple{#1}}}
\newcommand{\concrete}[1]{\ensuremath{\textit{concrete}(#1)}}

\newcommand{\dom}{\ensuremath{\textit{dom}}}
\newcommand{\noselfdom}{\ensuremath{\textit{noSelfDom}}}
\newcommand{\arrow}{\ensuremath{\textit{arrow}}}
\newcommand{\FV}{\ensuremath{\textit{FV}}}

% indented code block
\newenvironment{ttquote}%
{\begin{quote}\begin{alltt}}
{\end{alltt}\end{quote}}

\newcommand{\verythinmathspace}{\mskip0.5\thinmuskip}

\newcommand{\emptyBecomesThin}[1]{{\def\TEMPA{#1}\def\TEMPB{}\ifx\TEMPA\TEMPB \,\else #1 \fi}}
\newcommand{\emptyBecomesVeryThin}[1]{{\def\TEMPA{#1}\def\TEMPB{}\ifx\TEMPA\TEMPB \verythinmathspace \else #1 \fi}}

\newcommand{\set}[1]{\ensuremath{\{\emptyBecomesThin{#1}\}}}
\newcommand{\setb}[1]{\ensuremath{\{\bar{#1}\}}}
\newcommand{\bigset}[1]{\ensuremath{\big\{\emptyBecomesThin{#1}\big\}}}
\newcommand{\bigsetb}[1]{\ensuremath{\big\{\bar{#1}\big\}}}
\newcommand{\LRset}[1]{\ensuremath{\left\{\emptyBecomesThin{#1}\right\}}}
\newcommand{\LRsetb}[1]{\ensuremath{\left\{\bar{#1}\right\}}}

\newcommand{\Bigsetcomp}[2]{\ensuremath{\Bigl\{\,{#1}\;\Bigm|\;{#2}\,\Bigr\}}}
\newcommand{\LRsetcomp}[2]{\ensuremath{\left\{\,{#1}\;\left|\;{#2}\,\right\}\right.}}

\newcommand{\myemptyset}{\set{\,}}

\newcommand{\biglistof}[1]{\ensuremath{\bigl\langle\emptyBecomesVeryThin{#1}\bigr\rangle}}
\newcommand{\biglistofb}[1]{\ensuremath{\bigl\langle\rep{#1}\bigr\rangle}}
\newcommand{\Biglistof}[1]{\ensuremath{\Big\langle\emptyBecomesVeryThin{#1}\Big\rangle}}
\newcommand{\Biglistofb}[1]{\ensuremath{\Big\langle\rep{#1}\Big\rangle}}
\newcommand{\LRlistof}[1]{\ensuremath{\left\langle\emptyBecomesVeryThin{#1}\right\rangle}}
\newcommand{\LRlistofb}[1]{\ensuremath{\left\langle\rep{#1}\right\rangle}}
\newcommand{\emptylist}{\listof{\verythinmathspace}}

% tree structure for match
\newcommand{\emptytree}{\ensuremath{\bigl(\bigr)}}
\newcommand{\treel}[1]{\ensuremath{\bigl(#1,\set{}\bigr)}}
\newcommand{\tree}[2]{\ensuremath{\bigl(#1,\set{#2}\bigr)}}
\newcommand{\treeb}[2]{\ensuremath{\bigl(#1,\setb{#2}\bigr)}}
\newcommand{\node}[3]{\ensuremath{(#1\;#2\;#3)}}

\renewcommand{\root}{\textit{root}}
\newcommand{\children}{\textit{children}}
\newcommand{\nodes}{\textit{nodes}}

% put in oxford brackets
\newcommand{\ob}[1]{\ensuremath{\llbracket\verythinmathspace \emptyBecomesVeryThin{#1} \verythinmathspace\rrbracket}}
\newcommand{\bigob}[1]{\ensuremath{\big\llbracket\verythinmathspace \emptyBecomesVeryThin{#1} \verythinmathspace\big\rrbracket}}
\newcommand{\Bigob}[1]{\ensuremath{\Big\llbracket\verythinmathspace \emptyBecomesVeryThin{#1} \verythinmathspace\Big\rrbracket}}
% put in oxford brackets and an overbar
\newcommand{\obb}[1]{\ensuremath{\llbracket\verythinmathspace \bar{#1} \verythinmathspace\rrbracket}}
\newcommand{\bigobb}[1]{\ensuremath{\big\llbracket\verythinmathspace \bar{#1} \verythinmathspace\big\rrbracket}}
\newcommand{\Bigobb}[1]{\ensuremath{\Big\llbracket\verythinmathspace \bar{#1} \verythinmathspace\Big\rrbracket}}
% make a type param bound
% \newcommand{\bd}[1]{\ensuremath{\{{#1}\}}}
% \newcommand{\bigbd}[1]{\ensuremath{\bigl\{{#1}\bigr\}}}
\newcommand{\bd}[1]{{\def\TEMPA{#1}\def\TEMPB{\dontcares}\ifx\TEMPA\TEMPB \dontcare \else \ensuremath{\{{#1}\}}\fi}}
\newcommand{\bigbd}[1]{{\def\TEMPA{#1}\def\TEMPB{\dontcares}\ifx\TEMPA\TEMPB \dontcare \else \ensuremath{\bigl\{{#1}\bigr\}}\fi}}
\newcommand{\bdb}[1]{\ensuremath{\{\bar{#1}\}}}
\newcommand{\bigbdb}[1]{\ensuremath{\bigl\{\bar{#1}\bigr\}}}
% syntactic definition
\newcommand{\syndef}{\ensuremath{\overset{\mathrm{def}}{=}}}
% make a substitution
\newcommand{\subst}[2]{\ensuremath{[#1/#2]}}
\newcommand{\substb}[2]{\ensuremath{[\bar{#1/#2}]}}
\newcommand{\substbx}[4]{\ensuremath{[\bar{#1/#2},#3/#4]}}
\newcommand{\substxb}[4]{\ensuremath{[#1/#2,\bar{#3/#4}]}}
\newcommand{\substxx}[4]{\ensuremath{[#1/#2,#3/#4]}}
\newcommand{\substbb}[4]{\ensuremath{[\bar{#1/#2},\bar{#3/#4}]}}
\newcommand{\bigsubst}[2]{\ensuremath{\bigl[#1/#2\bigr]}}
\newcommand{\bigsubstb}[2]{\ensuremath{\bigl[\bar{#1/#2}\bigr]}}
\newcommand{\bigsubstbx}[4]{\ensuremath{\bigl[\bar{#1/#2},#3/#4\bigr]}}
\newcommand{\bigsubstxb}[4]{\ensuremath{\bigl[#1/#2,\bar{#3/#4}\bigr]}}
\newcommand{\bigsubstxx}[4]{\ensuremath{\bigl[#1/#2,#3/#4\bigr]}}
\newcommand{\Bigsubst}[2]{\ensuremath{\Bigl[#1/#2\Bigr]}}
\newcommand{\Bigsubstb}[2]{\ensuremath{\Bigl[\bar{#1/#2}\Bigr]}}
\newcommand{\Bigsubstbx}[4]{\ensuremath{\Bigl[\bar{#1/#2},#3/#4\Bigr]}}
\newcommand{\Bigsubstxb}[4]{\ensuremath{\Bigl[#1/#2,\bar{#3/#4}\Bigr]}}
\newcommand{\Bigsubstxx}[4]{\ensuremath{\Bigl[#1/#2,#3/#4\Bigr]}}
\newcommand{\emptysubst}{\ensuremath{[\verythinmathspace]}}
% type parameter list with bounds and oxford brackets
\newcommand{\tplist}[2]{\ensuremath{\ob{\bds{#1}{#2}}}}
% monomorphic fn decl
\newcommand{\decl}[3]{\mbox{\ensuremath{{#1}\,{#2}\!:\!{#3}}}}
% a generic function declaration 
\newcommand{\declg}[5]{\mbox{\ensuremath{#1 \tplist{#2}{#3}\, #4\!:\!#5}}}
\newcommand{\hdeclg}[4]{\mbox{\ensuremath{#1 \ob{#2}\, #3\!:\!#4}}}
% type context
\newcommand{\T}{\ensuremath{\mathcal{T}}}
% context application
\newcommand{\context}[1]{\,\setlength{\fboxsep}{1pt}\fbox{\ensuremath{\strut#1}}}
% empty context
\newcommand{\emptycontext}{\ensuremath{\Box}}
% class table extension
\newcommand{\ctext}{\ensuremath{\supseteq}}
% a declaration set D
\newcommand{\D}{\ensuremath{\mathcal{D}}}
% a declaration set restricted to a function namemeet	
\newcommand{\Df}[1][f]{\D_{\!#1}}
% existential type
\newcommand{\existstype}[2]{\ensuremath{\exists\ob{#1}{#2}}}
\newcommand{\existstypeb}[2]{\ensuremath{\exists\obb{#1}{#2}}}
\newcommand{\existstypebig}[2]{\ensuremath{\exists\bigob{#1}{#2}}}
\newcommand{\existstypebigb}[2]{\ensuremath{\exists\bigobb{#1}{#2}}}
\newcommand{\existstypeBig}[2]{\ensuremath{\exists\Bigob{#1}{#2}}}
\newcommand{\existstypeBigb}[2]{\ensuremath{\exists\Bigobb{#1}{#2}}}
% universal type
\newcommand{\foralltype}[2]{\ensuremath{\forall\ob{#1}{#2}}}
\newcommand{\foralltypeb}[2]{\ensuremath{\forall\obb{#1}{#2}}}
\newcommand{\foralltypebig}[2]{\ensuremath{\forall\bigob{#1}{#2}}}
\newcommand{\foralltypebigb}[2]{\ensuremath{\forall\bigobb{#1}{#2}}}
\newcommand{\foralltypeBig}[2]{\ensuremath{\forall\Bigob{#1}{#2}}}
\newcommand{\foralltypeBigb}[2]{\ensuremath{\forall\Bigobb{#1}{#2}}}
% reduced existential type
\newcommand{\reduce}[1]{\ensuremath{{#1}_r}}

% binding of type parameters with (name, upper bounds)
\newcommand{\bindup}[2]{#1 \extends \bd{#2}}
\newcommand{\bigbindup}[2]{#1 \extends \bigbd{#2}}
% binding of type parameters with (lower bounds, name, upper bounds)
\newcommand{\lobindup}[3]{\bd{#1} \extends #2 \extends \bd{#3}}
\newcommand{\lobinduptight}[3]{\bd{#1} \mathord{\extends} #2 \mathord{\extends} \bd{#3}}
\newcommand{\biglobindup}[3]{\bigbd{#1} \extends #2 \extends \bigbd{#3}}

% trait declaration (name, type parameters, extends, methods)
\newcommand{\traitdecl}[4]{{\def\TEMPA{#1}\def\TEMPB{#2}\def\TEMPC{#3}\def\TEMPD{#4}\def\TEMPX{\dontcare}\def\TEMPZ{\dontcares}\relax
    \ifx\TEMPD\TEMPZ
        \ifx\TEMPC\TEMPZ
            \ifx\TEMPB\TEMPZ
                \ifx\TEMPA\TEMPX
                    \kwd{trait} \; \dontcares \; \kwd{end}
                \else
                    \kwd{trait} \; #1\bigob{#2} \; \dontcares \; \kwd{end}
                \fi
            \else
                \kwd{trait} \; #1\bigob{#2} \; \dontcares \; \kwd{end}
            \fi
        \else
            \kwd{trait} \; #1\bigob{#2} \extends \bigbd{#3} \; #4 \; \kwd{end}
        \fi
    \else
        \kwd{trait} \; #1\bigob{#2} \extends \bigbd{#3} \; #4 \; \kwd{end}
    \fi}}
% object declaration (name, type parameters, value parameters, extends, methods)
\newcommand{\objectdecl}[5]{{\def\TEMPA{#1}\def\TEMPB{#2}\def\TEMPC{#3}\def\TEMPD{#4}\def\TEMPE{#5}\def\TEMPX{\dontcare}\def\TEMPZ{\dontcares}\relax
    \ifx\TEMPE\TEMPZ
        \ifx\TEMPD\TEMPZ
            \ifx\TEMPC\TEMPZ
                \ifx\TEMPB\TEMPZ
                    \ifx\TEMPA\TEMPX
                        \kwd{object} \; \dontcares \; \kwd{end}
                    \else
                        \kwd{object} \; #1\bigob{#2} \; \dontcares \; \kwd{end}
                    \fi
                \else
                    \kwd{object} \; #1\bigob{#2} \; \dontcares \; \kwd{end}
                \fi
            \else
                \kwd{object} \; #1\bigob{#2}\tuple{#3} \; \dontcares \; \kwd{end}
            \fi
        \else
            \kwd{object} \; #1\bigob{#2}\tuple{#3} \extends \bigbd{#4} \; #5 \; \kwd{end}
        \fi
    \else
        \kwd{object} \; #1\bigob{#2}\tuple{#3} \extends \bigbd{#4} \; #5 \; \kwd{end}
    \fi}}
% method definition (name, type parameters, value parameters, result type, body)
\newcommand{\methoddef}[5]{#1\ob{#2}\tuple{#3}\COLON#4 \defneq #5}

%%%%% Any and Bottom %%%%

\newcommand{\Any}{\TYP{Any}\xspace}
\newcommand{\Object}{\TYP{Object}\xspace}
\newcommand{\Bottom}{\TYP{Bottom}\xspace}

\newcommand{\bindco}{\hbox{\tt +}}
\newcommand{\bindcontra}{\hbox{\tt -}}
\newcommand{\bindinv}{\hbox{\tt =}}

\newcommand{\useco}{\hbox{\setbox0\hbox{$\bigcirc$}\hbox to \wd0{\hbox to 0pt{\ht0=0pt\raise0.1ex\box0\hss}\hss\tt +\hss}}}
\newcommand{\usecontra}{\hbox{\setbox0\hbox{$\bigcirc$}\hbox to \wd0{\hbox to 0pt{\ht0=0pt\raise0.1ex\box0\hss}\hss\tt -\hss}}}
\newcommand{\useinv}{\hbox{\setbox0\hbox{$\bigcirc$}\hbox to \wd0{\hbox to 0pt{\ht0=0pt\raise0.1ex\box0\hss}\hss\tt =\hss}}}

\newcommand{\FALSE}{\textrm{false}}
\newcommand{\TRUE}{\textrm{true}}

\newcommand{\NONE}{\bullet}

\newcommand{\eqred}{\overset{\equiv}{\longrightarrow}}

%%%%%%% JUDGMENTS %%%%%%%%

%%% NEW SYNTACTIC JUDGMENT
%%%\newcommand{\newjudge}[2]{\fbox{\textbf{#1:} \quad \ensuremath{#2}}}
% \newcommand{\newjudge}[2]{\hbox{\textbf{#1:}\hskip0.5em\fbox{\ensuremath{#2}}}}
% \newcommand{\newjudgetwo}[2]{\textbf{#1:}\vskip0em\hfill\fbox{\ensuremath{#2}}}

\newcommand{\newjudge}[2]{{\setbox0\hbox{\textbf{#1:}}\setbox1\hbox{\fbox{\ensuremath{#2}}}\setbox2\hbox{\copy0\hskip0.5em\copy1}\relax
  \ifdim \wd2 > \linewidth
     \hbox{\box0\vrule width 0pt depth 0.4em}\relax
     \hbox to \linewidth{\hss\box1}\relax
  \else
     \box2
  \fi}}

%\newcommand{\newjudgetwo}{\newjudge}

% non constrained judgements
\newcommand{\jgtemplate}[4][\Delta]{\ensuremath{\whennonempty{#1}{{#1}\,}\turnstile\,{#2}\;{#3}\;{#4}}}
\newcommand{\jgbtemplate}[4][\Delta]{\ensuremath{\rep{\whennonempty{#1}{\same{#1}\,}\turnstile\,{#2}\;{#3}\;{#4}}}}
\newcommand{\jgbbtemplate}[4][\Delta]{\ensuremath{\rep{\rep{\whennonempty{#1}{\same{\same{#1}}\,}\turnstile\,{#2}\;{#3}\;{#4}}}}}
\newcommand{\jgbARROWtemplate}[4][\Delta]{\ensuremath{\whennonempty{#1}{{#1}\,}\turnstile\,\bar{\overcirclerightarrow{{#2}\;{#3}\;{#4}}}}}
\newcommand{\jgtemplatetight}[4][\Delta]{\ensuremath{\whennonempty{#1}{{#1}\verythinmathspace}\turnstile\verythinmathspace{#2}\,{#3}\,{#4}}}
%\newcommand{\jgbtemplatetight}[4][\Delta]{\ensuremath{\whennonempty{#1}{{#1}\verythinmathspace}\turnstile\verythinmathspace\bar{{#2}\,{#3}\,{#4}}}}
\newcommand{\jgbtemplatetight}[4][\Delta]{\ensuremath{\rep{\whennonempty{#1}{\same{#1}\verythinmathspace}\turnstile\verythinmathspace{#2}\,{#3}\,{#4}}}}
\newcommand{\jgTHREEtemplate}[8]{\ensuremath{\whennonempty{#1}{{#1}\,}\turnstile\,{#2}\;{#3}\;{#4}, {#5}\;{#3}\;{#6}, {#7}\;{#3}\;{#8}}}

\newcommand{\jgshorttemplate}[3][\Delta]{\ensuremath{\whennonempty{#1}{{#1}\,}\turnstile\,{#2}\;{#3}}}
\newcommand{\jgbshorttemplate}[3][\Delta]{\ensuremath{\rep{\whennonempty{#1}{\same{#1}\,}\turnstile\,{#2}\;{#3}}}}
\newcommand{\jgbbshorttemplate}[3][\Delta]{\ensuremath{\rep{\rep{\whennonempty{#1}{\same{\same{#1}}\,}\turnstile\,{#2}\;{#3}}}}}
\newcommand{\jgshorttemplatetight}[3][\Delta]{\ensuremath{\whennonempty{#1}{{#1}\verythinmathspace}\turnstile\verythinmathspace{#2}\;{#3}}}
\newcommand{\jgbshorttemplatetight}[3][\Delta]{\ensuremath{\whennonempty{#1}{{#1}\verythinmathspace}\turnstile\verythinmathspace\bar{{#2}\;{#3}}}}

% ground subtyping
\newcommand{\jgsub}[3][\Delta]{\jgtemplate[#1]{#2}{\subtypeof}{#3}}
\newcommand{\jgnequiv}[3][\Delta]{\jgtemplate[#1]{#2}{\not \equiv}{#3}}

% evaluation
\newcommand{\denv}{\Psi}
\newcommand{\jevalstep}[3][\denv]{\jgtemplate[#1]{#2}{\longrightarrow}{#3}}
\newcommand{\jevalsteptight}[3][\Delta]{\jgtemplatetight[#1]{#2}{\longrightarrow}{#3}}
\newcommand{\jevalstepTWO}[5][\Delta]{\ensuremath{\whennonempty{#1}{{#1}\,}\turnstile\,{#2}\;{\longrightarrow}\;{#3}, {#4}\;{\longrightarrow}\;{#5}}}
\newcommand{\jbevalstep}[3][\Delta]{\jgbtemplate[#1]{#2}{\longrightarrow}{#3}}
\newcommand{\jevalstar}[3][\Delta]{\jgtemplate[#1]{#2}{\longrightarrow^*}{#3}}

% typing
\newcommand{\jtype}[3][\Delta;\Gamma]{\jgtemplate[#1]{#2}{\mathrel{:}}{#3}}
\newcommand{\jbtype}[3][\Delta;\Gamma]{\jgbtemplate[#1]{#2}{\mathrel{:}}{#3}}
\newcommand{\jtypetight}[3][\Delta;\Gamma]{\jgtemplatetight[#1]{#2}{\mathrel{:}}{#3}}
\newcommand{\jbtypetight}[3][\Delta;\Gamma]{\jgntemplatetight[#1]{#2}{\mathrel{:}}{#3}}

%variance
\newcommand{\jcertainlyvaries}[4][\Delta]{\jgtemplate[#1]{#2}{#3}{#4}}
\newcommand{\jbcertainlyvaries}[4][\Delta]{\jgbtemplate[#1]{#2}{#3}{#4}}

%subtyping
\newcommand{\jsubtype}[3][\Delta]{\jgtemplate[#1]{#2}{\subtypeof}{#3}}
\newcommand{\jbsubtype}[3][\Delta]{\jgbtemplate[#1]{#2}{\subtypeof}{#3}}
\newcommand{\jbbsubtype}[3][\Delta]{\jgbbtemplate[#1]{#2}{\subtypeof}{#3}}
\newcommand{\jnotsubtype}[3][\Delta]{\jgtemplate[#1]{#2}{\not\subtypeof}{#3}}
\newcommand{\jbnotsubtype}[3][\Delta]{\jgbtemplate[#1]{#2}{\not\subtypeof}{#3}}
\newcommand{\jbbnotsubtype}[3][\Delta]{\jgbbtemplate[#1]{#2}{\not\subtypeof}{#3}}
\newcommand{\jexctype}[3][\Delta]{\jgtemplate[#1]{#2}{\exc}{#3}}
\newcommand{\jbexctype}[3][\Delta]{\jgbtemplate[#1]{#2}{\exc}{#3}}
\newcommand{\jnotexctype}[3][\Delta]{\jgtemplate[#1]{#2}{\nexc}{#3}}
\newcommand{\jequivtype}[3][\Delta]{\jgtemplate[#1]{#2}{\equiv}{#3}}
\newcommand{\jbequivtype}[3][\Delta]{\jgbtemplate[#1]{#2}{\equiv}{#3}}
\newcommand{\jnotequivtype}[3][\Delta]{\jgtemplate[#1]{#2}{\not\equiv}{#3}}

% well-formed types
\newcommand{\jwftype}[2][\Delta]{\jgshorttemplate[#1]{#2}{\mathsf{ok}}}
\newcommand{\jbwftype}[2][\Delta]{\jgbshorttemplate[#1]{#2}{\mathsf{ok}}}
\newcommand{\jbbwftype}[2][\Delta]{\jgbbshorttemplate[#1]{#2}{\mathsf{ok}}}
\newcommand{\jwftypetight}[2][\Delta]{\jgshorttemplatetight[#1]{#2}{\mathsf{ok}}}
\newcommand{\jbwftypetight}[2][\Delta]{\jgbshorttemplatetight[#1]{#2}{\mathsf{ok}}}

\newcommand{\jwfdecl}[2][\Delta]{\jgshorttemplate[#1]{#2}{\mathsf{ok}}}
\newcommand{\jbwfdecl}[2][\Delta]{\jgbshorttemplate[#1]{#2}{\mathsf{ok}}}

\newcommand{\jwfmeth}[2][\Delta;\Gamma;\bar{V\;P}]{\jgshorttemplate[#1]{#2}{\mathsf{ok}}}
\newcommand{\jbwfmeth}[2][\Delta;\Gamma;\bar{V\;P}]{\jgbshorttemplate[#1]{#2}{\mathsf{ok}}}

% Meet and return type rules
\newcommand{\jgruletemplate}[5][\Delta]{{#1}\,\turnstile\,{#2}\;{#3}\;{#4}\;{#5}}
\newcommand{\jgbruletemplate}[5][\Delta]{{#1}\,\turnstile\,\bar{{#2}\;{#3}\;{#4}\;{#5}}}
\newcommand{\jgbbruletemplate}[5][\Delta]{{#1}\,\turnstile\,\bar{\bar{{#2}\;{#3}\;{#4}\;{#5}}}}
\newcommand{\jgbARROWruletemplate}[5][\Delta]{{#1}\,\turnstile\,\bar{\overcirclerightarrow{{#2}\;{#3}\;{#4}\;{#5}}}}
\newcommand{\jmethodmeetrule}[4][\Delta]{{#1}\,\turnstile\,{#2}\;\mathsf{meet}\;{#3}\;\mathsf{wrt}\;{#4}\;\mathsf{ok}}
%\newcommand{\jmethodmeetrule}[3][\Delta]{\jgruletemplate[#1]{#2}{\mathsf{meet}}{#3}{\mathsf{ok}}}
\newcommand{\jbmethodmeetrule}[3][\Delta]{\jgbruletemplate[#1]{#2}{\mathsf{meet}}{#3}{\mathsf{ok}}}
\newcommand{\jbbmethodmeetrule}[3][\Delta]{\jgbbruletemplate[#1]{#2}{\mathsf{meet}}{#3}{\mathsf{ok}}}
\newcommand{\jbARROWmethodmeetrule}[3][\Delta]{\jgbARROWruletemplate[#1]{#2}{\mathsf{meet}}{#3}{\mathsf{ok}}}
\newcommand{\jreturntyperule}[3][\Delta]{\jgruletemplate[#1]{#2}{\mathsf{return\;type\;wrt}}{#3}{\mathsf{ok}}}
\newcommand{\jbreturntyperule}[3][\Delta]{\jgbruletemplate[#1]{#2}{\mathsf{return\;type\;wrt}}{#3}{\mathsf{ok}}}
\newcommand{\jbbreturntyperule}[3][\Delta]{\jgbbruletemplate[#1]{#2}{\mathsf{return\;type\;wrt}}{#3}{\mathsf{ok}}}
\newcommand{\jbARROWreturntyperule}[3][\Delta]{\jgbARROWruletemplate[#1]{#2}{\mathsf{return\;type\;wrt}}{#3}{\mathsf{ok}}}

\newcommand{\jgfunnyruletemplate}[5][\Delta]{{#1}\,\turnstile\,{#2}\;\mathsf{and}\;{#3}\;{#4}\;{#5}}
\newcommand{\jgbfunnyruletemplate}[5][\Delta]{{#1}\,\turnstile\,\bar{{#2}\;\mathsf{and}\;{#3}\;{#4}\;{#5}}}
\newcommand{\jgbbfunnyruletemplate}[5][\Delta]{{#1}\,\turnstile\,\bar{\bar{{#2}\;\mathsf{and}\;{#3}\;{#4}\;{#5}}}}
\newcommand{\jgbARROWfunnyruletemplate}[5][\Delta]{{#1}\,\turnstile\,\bar{\overcirclerightarrow{{#2}\;\mathsf{and}\;{#3}\;{#4}\;{#5}}}}
\newcommand{\jancestormeetrule}[3][\Delta;\bar{J}]{\jgfunnyruletemplate[#1]{#2}{#3}{\mathsf{ancestors}}{\mathsf{ok}}}
\newcommand{\jbancestormeetrule}[3][\Delta;\bar{J}]{\jgbfunnyruletemplate[#1]{#2}{#3}{\mathsf{ancestors}}{\mathsf{ok}}}
\newcommand{\jbbancestormeetrule}[3][\Delta;\bar{J}]{\jgbbfunnyruletemplate[#1]{#2}{#3}{\mathsf{ancestors}}{\mathsf{ok}}}
\newcommand{\jbARROWancestormeetrule}[3][\Delta;\bar{\chi}]{\jgbARROWfunnyruletemplate[#1]{#2}{#3}{\mathsf{ancestors}}{\mathsf{ok}}}

\newcommand{\jnoduplicatesrule}[3][\Delta]{\jgtemplate[#1]{#2}{\mathsf{not}\;\mathsf{duplicate}\;\mathsf{of}}{#3}}
\newcommand{\jbnoduplicatesrule}[3][\Delta]{\jgbtemplate[#1]{#2}{\mathsf{not}\;\mathsf{duplicate}\;\mathsf{of}}{#3}}
\newcommand{\jbbnoduplicatesrule}[3][\Delta]{\jgbbtemplate[#1]{#2}{\mathsf{not}\;\mathsf{duplicate}\;\mathsf{of}}{#3}}
\newcommand{\jbARROWnoduplicatesrule}[3][\Delta]{\jgbARROWtemplate[#1]{#2}{\mathsf{not}\;\mathsf{duplicate}\;\mathsf{of}}{#3}}

\newcommand{\jcoveredrule}[3][\Delta]{\jgtemplate[#1]{#2}{\mathsf{covered}\;\mathsf{by}}{#3}}
\newcommand{\jbcoveredrule}[3][\Delta]{\jgbtemplate[#1]{#2}{\mathsf{covered}\;\mathsf{by}}{#3}}
\newcommand{\jbbcoveredrule}[3][\Delta]{\jgbbtemplate[#1]{#2}{\mathsf{covered}\;\mathsf{by}}{#3}}
\newcommand{\jbARROWcoveredrule}[3][\Delta]{\jgbARROWtemplate[#1]{#2}{\mathsf{covered}\;\mathsf{by}}{#3}}

% subtyping on quantified types
\newcommand{\jle}[3][\Delta]{\jgtemplate[#1]{#2}{\quantsubtypeof}{#3}}
\newcommand{\jleusing}[4][\Delta]{\jgtemplate[#1]{#2}{\quantsubtypeof}{#3}\;\mathsf{using}\;{#4}}
\newcommand{\jleusingtight}[4][\Delta]{\jgtemplatetight[#1]{#2}{\quantsubtypeof}{#3}\;\mathsf{using}\;{#4}}
\newcommand{\jleinner}[3][\Delta]{\jgtemplate[#1]{#2}{\leinner}{#3}}
\newcommand{\jleinnerusing}[4][\Delta]{\jgtemplate[#1]{#2}{\leinner}{#3}\;\mathsf{using}\;{#4}}
\newcommand{\jleinnerusingtight}[4][\Delta]{\jgtemplatetight[#1]{#2}{\leinner}{#3}\;\mathsf{using}\;{#4}}
\newcommand{\jexcludes}[3][\Delta]{\jgtemplate[#1]{#2}{\existsexc}{#3}}
\newcommand{\jequivalent}[3][\Delta]{\jgtemplate[#1]{#2}{\equiv}{#3}}

\newcommand{\curriedmorespecific}[1][\Delta]{\textit{moreSpecific}\tuple{#1}}
\newcommand{\morespecific}[3][\Delta]{\curriedmorespecific[#1]{\tuple{{#2},{#3}}}}
\newcommand{\jmorespecific}[3][\Delta]{\jgtemplate[#1]{\dom(#2)}{\morespec}{\dom(#3)}}
\newcommand{\jmorespecificusing}[4][\Delta]{\jgtemplate[#1]{\dom(#2)}{\morespec}{\dom(#3)}\;\mathsf{using}\;{#4}}


% constrained judgments
\newcommand{\jgconstrtemplate}[5][\Delta]{\ensuremath{\whennonempty{#1}{{#1}\,}\turnstile\,{#2}\;{#3}\;{#4}\,\Leftarrow\,{#5}}}
\newcommand{\jgbconstrtemplate}[5][\Delta]{\ensuremath{\whennonempty{#1}{{#1}\,}\turnstile\,\bar{{#2}\;{#3}\;{#4}\,\Leftarrow\,{#5}}}}
\newcommand{\jgbbconstrtemplate}[5][\Delta]{\ensuremath{\whennonempty{#1}{{#1}\,}\turnstile\,\bar{\bar{{#2}\;{#3}\;{#4}\,\Leftarrow\,{#5}}}}}
% ground subtyping with constraints
\newcommand{\jsub}[4][\Delta]{\jgconstrtemplate[#1]{#2}{\subtypeof}{#3}{#4}}
\newcommand{\jbsub}[4][\Delta]{\jgbconstrtemplate[#1]{#2}{\subtypeof}{#3}{#4}}
% not subtype
\newcommand{\jnsub}[4][\Delta]{\jgconstrtemplate[#1]{#2}{\not \subtypeof}{#3}{#4}}
\newcommand{\jbnsub}[4][\Delta]{\jgbconstrtemplate[#1]{#2}{\not \subtypeof}{#3}{#4}}
% type exclusion
\newcommand{\jexc}[4][\Delta]{\jgconstrtemplate[#1]{#2}{\exc}{#3}{#4}}
\newcommand{\jbexc}[4][\Delta]{\jgbconstrtemplate[#1]{#2}{\exc}{#3}{#4}}
\newcommand{\jmexc}[4][\Delta]{\jgconstrtemplate[#1]{#2}{\bexcm}{#3}{#4}}
% type non-exclusion
\newcommand{\jnexc}[4][\Delta]{\jgconstrtemplate[#1]{#2}{\nexc}{#3}{#4}}
% equivalence
\newcommand{\jequiv}[4][\Delta]{\jgconstrtemplate[#1]{#2}{\equiv}{#3}{#4}}
% nonequivalence
\newcommand{\jnequiv}[4][\Delta]{\jgconstrtemplate[#1]{#2}{\not\equiv}{#3}{#4}}
\newcommand{\jbnequiv}[4][\Delta]{\jgbconstrtemplate[#1]{#2}{\not\equiv}{#3}{#4}}

% ground subtyping with constraints
\newcommand{\jvaries}[5][\Delta]{\jgconstrtemplate[#1]{#2}{#3}{#4}{#5}}
\newcommand{\jbvaries}[5][\Delta]{\jgbconstrtemplate[#1]{#2}{#3}{#4}{#5}}

\newcommand{\jcomp}[5][\Delta]{\ensuremath{\whennonempty{#1}{{#1}\,}\turnstile\,{#2}\;\mathsf{and}\;{#3}\;\mathsf{compare}\;{#4}\,\Leftarrow\,{#5}}}
%\newcommand{\jcomp}[5][\Delta]{\jgconstrtemplate[#1]{#2 \;\mathsf{and}\;#3}\;{\mathsf{compare}\;#4}\;{#5}}
%\newcommand{\jcompare}[4][\Delta]{\ensuremath{\whennonempty{#1}{{#1}\,}\turnstile\,{#2}\;\mathsf{and}\;{#3}\;\mathsf{compare}\;{#4}}}

% contrapositive judgements
\newcommand{\jgcontratemplate}[5][\Delta]{\ensuremath{\whennonempty{#1}{{#1}\,}\turnstile\,{#2}\;{#3}\;{#4}\,\Rightarrow\,{#5}}}
% ground subtyping with constraints
\newcommand{\jcsub}[4][\Delta]{\jgcontratemplate[#1]{#2}{\subtypeof}{#3}{#4}}
% not subtype
\newcommand{\jcnsub}[4][\Delta]{\jgcontratemplate[#1]{#2}{\not \subtypeof}{#3}{#4}}
% type exclusion
\newcommand{\jcexc}[4][\Delta]{\jgcontratemplate[#1]{#2}{\exc}{#3}{#4}}
% type non-exclusion
\newcommand{\jcnexc}[4][\Delta]{\jgcontratemplate[#1]{#2}{\nexc}{#3}{#4}}
% equivalence
\newcommand{\jcequiv}[4][\Delta]{\jgcontratemplate[#1]{#2}{\equiv}{#3}{#4}}
% nonequivalence
\newcommand{\jcnequiv}[4][\Delta]{\jgcontratemplate[#1]{#2}{\not\equiv}{#3}{#4}}


% applicability of a domain or fndecl to a type
\newcommand{\japp}[3][\Delta]{\jgtemplate[#1]{#2}{\ni}{#3}}
% specificity between fndecls
\newcommand{\jms}[3][\Delta]{\jgtemplate[#1]{#2}{\ms}{#3}}

% constraints
% convert a bound environment into a constraint
\newcommand{\toConstraint}[2]{\ensuremath{\textit{toConstraint}({#1})\,=\,{#2}}}
% convert a constraint into a bound environment
\newcommand{\toBounds}[2]{\ensuremath{\textit{toBounds}({#1})\,=\,{#2}}}

% apply substitution to constraint
\newcommand{\japply}[4][\Delta]{\ensuremath{\whennonempty{#1}{{#1}\,}\turnstile\,\textit{apply}({#2}, {#3})\,=\,{#4}}}
% solve constraint to get a substitution and the residual constraints
\newcommand{\junify}[4][\Delta]{\ensuremath{\whennonempty{#1}{{#1}\,}\turnstile\,\textit{unify}({#2})\,=\,({#3}, {#4})}}

% type reduction
\newcommand{\jtred}[3][\Delta]{\ensuremath{\whennonempty{#1}{{#1}\,}\turnstile\,{#2} \eqred {#3}}}
\newcommand{\jtreds}[4][\Delta]{\ensuremath{\whennonempty{#1}{{#1}\,}\turnstile\,{#2} \eqred {#3}\;\mathsf{using}\;{#4}}}

\newcommand{\jscope}[2][\Delta]{\ensuremath{\whennonempty{#1}{{#1}\,}\turnstile\,{#2}\;\mathsf{ok}}}
\newcommand{\jextremal}[4][\Delta]{\ensuremath{\whennonempty{#1}{{#1}\,}\turnstile\,{#2}\;\mathsf{used\ in}\;{#3}\sqsubseteq{#4}\;\mathsf{extremal}}}

%matching
\newcommand{\jmatch}[7][\denv]{\ensuremath{\whennonempty{#1}{{#1}\,}\turnstile\,{#2}\;\mathsf{match}\;{#3}\;{#4}\;{#5}\;\mathsf{bounds}\;{#6}\;\mathsf{and}\;{#7}}}

\newcommand{\jcompare}[4][\Delta]{\ensuremath{\whennonempty{#1}{{#1}\,}\turnstile\,{#2}\;\mathsf{and}\;{#3}\;\mathsf{compare}\;{#4}}}
\newcommand{\jdonotcompare}[4][\Delta]{\ensuremath{\whennonempty{#1}{{#1}\,}\turnstile\,{#2}\;\mathsf{and}\;{#3}\;\mathsf{do}\;\mathsf{not}\;\mathsf{compare}\;{#4}}}
\newcommand{\jrelate}[4][\Delta]{\ensuremath{\whennonempty{#1}{{#1}\,}\turnstile\,{#2}\;\mathsf{and}\;{#3}\;\mathsf{relate}\;{#4}}}
\newcommand{\jobeys}[5][\Delta]{\ensuremath{\whennonempty{#1}{{#1}\,}\turnstile\,{#2}\;\mathsf{on}\;{#3}\;\mathsf{obeys}\;{#4}\;\mathsf{and}\;{#5}}}

\newcommand{\jsolve}[5][\denv]{\ensuremath{\whennonempty{#1}{{#1}\,}\turnstile\,\mathsf{solve}\;{#2}\;\mathsf{bounds}\;{#3}\;\mathsf{and}\;{#4}\;\mathsf{using}\;{#5}}}

\newcommand{\jdispatch}[5][\denv]{\ensuremath{\whennonempty{#1}{{#1}\,}\turnstile\,{#2}\;\mathsf{and}\;{#3}\;\mathsf{dispatch}\;{#4}\;\mathsf{using}\;{#5}}}
\newcommand{\jfancydispatch}[5][\denv]{\ensuremath{\whennonempty{#1}{{#1}\,}\turnstile\,{#2}\;\mathsf{and}\;{#3}\;\mathsf{dispatch}\;{#4}\;\mathsf{using}\;{#5}}}

\newcommand{\jcontextvar}[3][\Delta]{\ensuremath{\whennonempty{#1}{{#1}\,}\turnstile\,{#2}\;\mathsf{variance}\;{#3}}}

\newcommand\SPANTWO[1]{\multicolumn{2}{@{}l@{\hskip3em}}{#1}}
\newcommand\spantwo[1]{\multicolumn{2}{@{}l@{}}{#1}}

% for tabularx environments to have a right-aligned, stretched col
\newcolumntype{R}{>{\raggedleft\arraybackslash}X}%

% \theorembodyfont{\rm}
% \newtheorem{proposition}{Proposition}
% \newtheorem{lemma}{Lemma}
% \newtheorem{theorem}{Theorem}
% Our proofs are more like proof sketches!! EricAllen 7/15/2011
% \newenvironment{proof}{\noindent \textbf{Proof:} }{\hfill $\Box$}
\newenvironment{psketch}{\noindent Proof sketch: }{\hfill $\Box$}
